% Modernised reading — fonds Alexandre Grothendieck, shelfmark 151, batch 3
% (pages 41-60).
%
% This is an INTERPRETATION, not a transcription. It re-reads the pages in
% today's notation and vocabulary, states the hypotheses the manuscript leaves
% implicit, and drops the critical apparatus so that the mathematics reads
% straight through. Everything the brackets and underlines carried — a doubtful
% reading, a notation he used differently, a missing passage — moves into a
% footnote.
%
% It is held to being mathematically correct as it stands, not merely faithful
% to the page. Where the manuscript is loose or mistaken, the true statement is
% what appears here, and the footnote says what the page has. In any dispute of
% substance the transcription governs: whoever cites Grothendieck must cite
% batch-03.fr.tex.
%
% Scope: the folder was read WHOLE before any of this was written — all four
% batches, pages 1-74, in one pass — and one file is written per batch, this
% being the third. Two places depend on that: the reading of the mixed tubes
% V^j_{i,k} on page 59, which is corrected here against the use the last batch
% makes of them; and the reason section 2's colimit is stated as strict rather
% than homotopical, which only the last batch's gluing theorem makes visible.
%
% Pass: Opus 5 (claude-opus-5), 2026-08-16 — first pass, unchecked against the
% pages by a human.
\documentclass[11pt,a4paper]{article}
% Le préambule commun à toutes les transcriptions du fonds Grothendieck.
%
% Il tient en une idée : **l'incertitude se compose, elle ne se résout pas.**
% Une transcription qui lisse un mot douteux a détruit la seule chose pour
% laquelle on la faisait — un lecteur doit pouvoir savoir, à chaque endroit,
% ce qui était sur le feuillet et ce qui a été ajouté. Les six macros qui
% suivent sont donc l'appareil critique, et non de la décoration.
%
% Les mêmes commandes sont reconnues par `scripts/render.mjs`, qui produit la
% vue de lecture du site. Une macro ajoutée ici doit l'être aussi là-bas,
% faute de quoi le rendu échouera bruyamment — ce qui est voulu : un rendu
% silencieusement faux est pire qu'un rendu absent.

\ProvidesPackage{grothendieck}[2026/08/08 Transcriptions du fonds Grothendieck]

\usepackage{amsmath,amssymb,amsthm}
\usepackage{mathrsfs}
\usepackage{stmaryrd}
% Les diagrammes commutatifs sont partout dans ce fonds : c'est la langue
% même de la géométrie algébrique de Grothendieck, et une transcription qui
% les rendrait en prose ne transcrirait rien.
\usepackage{tikz-cd}
% Français par défaut — c'est la langue des feuillets — mais l'anglais est
% chargé aussi : la traduction et le résumé sont des documents anglais, et la
% typographie française (espaces avant les deux-points) y serait une faute.
% Ces documents basculent par \selectlanguage{english} en tête de corps.
\usepackage[english,french]{babel}
\usepackage{fontspec}
\usepackage{geometry}
\geometry{a4paper,margin=25mm,marginparwidth=28mm}
\usepackage{marginnote}
\usepackage{xcolor}
% ulem, not soul. `soul`'s \st and \ul are built on a character-by-character
% scan that XeLaTeX's font machinery breaks: the document fails with an
% undefined control sequence rather than degrading. ulem does the same two jobs
% and is Unicode-engine safe. `normalem` stops it hijacking \emph, which the
% transcriptions use for Grothendieck's own emphasis.
\usepackage[normalem]{ulem}
\usepackage{hyperref}
\hypersetup{colorlinks=true,linkcolor=black,urlcolor=black,pdfborder={0 0 0}}

% --- Métadonnées du lot -----------------------------------------------------
% Reprises telles quelles par le script de rendu, qui les affiche en tête de
% la vue de lecture. Elles fixent ce que le fichier prétend couvrir : sans
% elles, une transcription flotte sans point d'attache dans un fonds de seize
% mille feuillets.

\newcommand{\folder}[1]{\gdef\grothFolder{#1}}
\newcommand{\batch}[1]{\gdef\grothBatch{#1}}
\newcommand{\leaves}[2]{\gdef\grothFirst{#1}\gdef\grothLast{#2}}
\newcommand{\foldertitle}[1]{\gdef\grothFoldertitle{#1}}
\gdef\grothFolder{?}\gdef\grothBatch{?}\gdef\grothFirst{?}\gdef\grothLast{?}\gdef\grothFoldertitle{}

% --- Le résumé --------------------------------------------------------------
%
% Il ouvre la lecture modernisée, sous le titre « Résumé », et il est composé
% en petit. Ce n'est pas un ornement : c'est le seul passage du document qui
% ne soit pas des mathématiques, et un lecteur doit pouvoir le reconnaître
% avant de l'avoir lu — puis le sauter s'il connaît déjà le sujet, ou s'y
% arrêter s'il ne le connaît pas. Le filet qui le suit marque la reprise du
% corps.

\newenvironment{resume}
  {\small\setlength{\parskip}{0.45em}}
  {\par\medskip\hrule\medskip}

% --- Mots-clés ---------------------------------------------------------------
%
% En anglais, en clôture du résumé de la lecture modernisée : le vocabulaire
% moderne sous lequel chercher ce que les feuillets construisent. Le manifeste
% du site (`npm run manifest`) les extrait pour étiqueter la cote entière —
% cette ligne est la source unique des tags, il n'y a pas de fichier à côté.
\newcommand{\keywords}[1]{\par\smallskip\noindent
  {\footnotesize\textsc{Keywords} --- \textit{#1}}\par}

% --- L'appareil critique ----------------------------------------------------

% Le numéro de feuillet, en marge. C'est l'ancre qui permet de régler un
% désaccord sur une formule en regardant *un* feuillet plutôt qu'une chemise
% de six cents. Le site s'en sert aussi pour tourner les pages du fac-similé
% quand on fait défiler la transcription.
\newcommand{\leaf}[1]{%
  \marginnote{\footnotesize\color{gray!70}\textsf{#1}}%
  \label{leaf:#1}\ignorespaces}

% Illisible : on ne devine pas. Un mot inventé qui se lit comme les autres est
% le pire résultat possible.
\newcommand{\ill}{\textcolor{red!60!black}{[\,\ldots\,]}}

% Lecture douteuse : le mot est donné, et signalé comme incertain.
\newcommand{\uncertain}[1]{\uline{#1}}

% Ajout de l'éditeur — une lettre manquante, un mot sous-entendu.
\newcommand{\add}[1]{\textcolor{blue!50!black}{[#1]}}

% Biffé par Grothendieck lui-même. Ce qu'il a rayé fait partie du document :
% une rature dit souvent où la pensée a changé de direction.
\newcommand{\struck}[1]{\sout{#1}}

% Note du transcripteur, sur l'état matériel du feuillet.
\newcommand{\note}[1]{\par\smallskip{\small\color{gray!60!black}\textit{[#1]}}\par\smallskip}

% Note marginale de Grothendieck, distincte du corps du texte.
\newcommand{\marginal}[1]{\par\smallskip\noindent\hspace{1em}%
  {\small\color{gray!60!black}#1}\par\smallskip}

% --- Titre ------------------------------------------------------------------

\AtBeginDocument{%
  \begin{center}
    {\large\bfseries Fonds Alexandre Grothendieck --- cote n\textsuperscript{o}~\grothFolder}\\[2pt]
    {\itshape\grothFoldertitle}\\[4pt]
    {\small feuillets \grothFirst--\grothLast\ (lot \grothBatch)}\\[2pt]
    {\footnotesize\color{gray!60!black}Transcription établie d'après le fac-similé numérisé
     par l'Université de Montpellier.\\
     Les lectures incertaines sont soulignées, les illisibles notées [\,\ldots\,],
     les ajouts entre crochets.}
  \end{center}
  \vspace{1em}}

\endinput


\folder{151}
\batch{3}
\pages{41}{60}
\dating{[à partir de 1981-à partir de 1990]}
\watermark{Édition de démonstration\\interprétation personnelle de l'œuvre}
\foldertitle{Topos stratifié (1981) : notes manuscrites (s.d.) — lecture modernisée}

\begin{document}

\section*{Résumé}

\begin{resume}
Le lot précédent laissait le texte au milieu d'une phrase, et devant un échec :
la sphère, découpée en un point et un plan, n'était pas retrouvée par la donnée
des groupes fondamentaux de ses morceaux. Ces vingt pages commencent par
réparer cela, puis quittent le cas de deux morceaux pour le cas général.

La réparation tient en trois pages et elle est jolie. Le groupe fondamental est
un invariant de dimension~1 ; ce qui manquait est un cran plus haut, et il y a
un objet qui l'habite : la \emph{gerbe}. Une gerbe est à un groupe abélien $A$
ce qu'un fibré en droites est à ses fonctions de transition, un cran plus haut ;
et sur un espace dont les morceaux sont des $K(\pi,1)$, se donner une gerbe
revient à se donner deux extensions de groupes et un isomorphisme entre ce
qu'elles induisent sur la jointure. Sur la sphère, le compte est immédiat : il
reste un automorphisme de l'extension triviale de $\mathbb{Z}$ par $A$, donc un
élément de $A$, et l'on retrouve $H^{2}(S^{2},A) \simeq A$ — « un magnifique
succès ! », écrit-il. Ce que le niveau des groupes fondamentaux perdait, le
niveau des gerbes le retrouve.

Le reste du lot construit le cadre général. Les morceaux ne sont plus deux mais
une famille $(X_{i})$ de parties fermées, indexée par un ensemble ordonné $I$,
la relation $i \leq j$ voulant dire $X_{i} \subset X_{j}$. À cette famille est
attaché un objet combinatoire : l'ensemble $X_{\Delta_{r}}$ des \emph{drapeaux},
c'est-à-dire des points $x$ munis d'une chaîne croissante $i_{0} \leq \cdots \leq
i_{r}$ d'indices dont le premier contient $x$. Ces ensembles s'organisent en un
espace simplicial, et la première proposition dit que l'espace tout entier se
retrouve à partir d'eux : $X$ est la limite inductive du système.

C'est vrai — et insuffisant, pour la raison que la sphère a déjà montrée. La
limite inductive en question est une limite d'espaces, pas de types d'homotopie,
et les $X_{i}$ sont fermés les uns dans les autres : entre deux morceaux
voisins, il n'y a rien qu'un bord, et un bord ne dit pas comment l'un s'approche
de l'autre. C'est exactement pourquoi la troisième section, qui s'ouvre à la page
56, remplace les fermés $X_{i}$ par les \emph{tubes} $\mathcal{V}_{i,j}$ : des
voisinages, et non des bords. Toute la fin du lot est la mise en place de ce
vocabulaire — la strate $X_{i}^{*}$, le tube $\mathcal{V}_{i,j}$, le tube épointé
$\mathcal{V}^{*}_{i,j}$, et une famille de tubes mixtes qui interpolent entre les
deux.

Deux passages valent d'être lus pour eux-mêmes. Le premier est une question, à
la page 53 : faut-il exiger que l'intersection de deux morceaux soit encore un
morceau ? Il répond non, en dessinant un cercle coupé en deux arcs par deux
points — deux morceaux connexes dont l'intersection ne l'est pas — et conclut :
« faut-il l'éliminer pour autant ? Il ne semble pas. La suite le montrera bien ! »
Le second est une remarque sur sa propre pratique, page 43 : la théorie
d'« intégration » des types d'homotopie qu'il appelle de ses vœux, et qu'il
voyait naguère comme un lien au grand yoga des relations entre structures
homotopiques et groupoïdes, lui paraît, réflexion faite, « relativement triviale,
et justiciable sans doute de techniques d'explicitation des plus standard ». Le
jugement est juste : c'est la théorie des colimites homotopiques, et elle
existait déjà.

\keywords{gerbe, poset-stratified space, simplicial space, homotopy colimit, conical stratification, link of a stratum}
\end{resume}

\pagerange{41}{43}
\section*{Les gerbes, ou ce que le niveau $\Pi_{1}$ laissait échapper}

La phrase interrompue au bas de la page 40 se termine ici, et elle énonce un
principe de recollement d'un cran plus haut que celui des systèmes locaux : la
donnée d'une gerbe $\mathcal{G}_{X}$ sur $X$ équivaut à celle d'une gerbe
$\mathcal{G}_{X^{*}}$ sur $X^{*}$, d'une gerbe $\mathcal{G}_{Y,X}$ sur le tube
$\mathcal{V}_{Y,X}$, et d'une équivalence entre leurs restrictions au tube
épointé :
\[
  \mathcal{G}_{Y,X}\bigl|\,\mathcal{V}^{*}_{Y,X} \;\simeq\;
  \mathcal{G}_{X^{*}}\bigl|\,\mathcal{V}^{*}_{Y,X}.
\]
C'est le même schéma de recollement que pour les faisceaux constructibles, avec
des gerbes à la place des systèmes locaux — et, comme il s'agit d'une
équivalence entre objets d'une $2$-catégorie et non d'une égalité, c'est
précisément là que loge l'information que la troncation en dimension~1 avait
perdue.

Sur un $K(\pi,1)$, une $A$-gerbe s'interprète canoniquement comme une extension
de $\pi$ par $A$.\footnote{Pour $A$ abélien et le lien donné, les $A$-gerbes sur
$K(\pi,1)$ sont classées par $H^{2}(\pi,A)$, qui classe aussi les extensions de
$\pi$ par $A$ ; l'énoncé du manuscrit est celui-là. Il suppose, comme il le dit,
les pièces connexes et de $\pi_{2}$ nul.} Se donner une gerbe sur $X$ revient
donc à se donner deux extensions — l'une de $\pi_{1}(Y)$ par $A$, l'autre de
$\pi_{1}(X^{*})$ par $A$ — et un isomorphisme entre les deux extensions de $D$
par $A$ qu'elles induisent, $D = \pi_{1}(\mathcal{V}^{*}_{Y,X})$ étant le groupe
de décomposition, muni de ses deux flèches vers $\pi_{1}(Y)$ et
$\pi_{1}(X^{*})$.

\subsection*{Le calcul sur la sphère}

Reprenons $X = S^{2}$, $Y = \{\infty\}$, $X^{*} = \mathbb{C}$, où $D \simeq
\mathbb{Z}$ et où les deux groupes fondamentaux $\pi_{1}(Y)$ et
$\pi_{1}(X^{*})$ sont triviaux. Les deux extensions sont alors triviales, et il
ne reste rien à choisir que l'isomorphisme entre les deux extensions triviales
de $\mathbb{Z}$ par $A$ qu'elles induisent — c'est-à-dire un automorphisme de
l'extension triviale, c'est-à-dire un élément de
$\mathrm{Hom}(\mathbb{Z},A) \simeq A$. D'où
\[
  H^{2}(S^{2}, A) \;\simeq\; A,
\]
la valeur correcte.\footnote{L'interprétation de $H^{2}(X,A)$ comme classifiant
les gerbes liées par $A$ est celle de Giraud (1971), à qui le manuscrit ne
renvoie pas ici. Ce que le calcul établit n'est pas la valeur de
$H^{2}(S^{2},A)$, qui est connue, mais que le schéma de recollement la donne :
c'est un test du programme, non un résultat sur la sphère.}

Ce succès vaut confirmation de l'intuition d'ensemble : il y a une théorie des
$\varinjlim$ de types d'homotopie, et le type d'homotopie d'une limite inductive
de topos est la limite inductive des types d'homotopie. Pour des espaces
stratifiés dont les pièces sont des $K(\pi,1)$ — l'exemple visé étant les espaces
de modules de Teichmüller, leurs compactifiés et leurs voisinages tubulaires —
le type d'homotopie doit s'exprimer en termes d'un système inductif convenable,
le plus souvent fini, des groupoïdes fondamentaux des pièces. C'est le programme
que sert tout le reste du dossier.\footnote{« Cette théorie d'intégration des
types d'homotopie […] est en fait relativement triviale, et justiciable sans
doute de techniques d'explicitation des plus standard », écrit-il, en ajoutant
qu'il se réserve d'y revenir. Le jugement est exact et l'outil existait : les
colimites homotopiques de Bousfield--Kan (1972), et le théorème de Thomason
(1979) identifiant la colimite homotopique d'un diagramme de catégories à la
construction de Grothendieck, sont l'un et l'autre antérieurs à ces pages.}

\pagerange{44}{49}
\section*{2. Stratifications globales : le formalisme $\tfrac{1}{2}$ simplicial, sans tubes}

On quitte le cas de deux morceaux. Le cadre est un espace topologique $X$ — plus
tard, un topos quelconque — et une famille de sous-espaces indexée par un
ensemble ordonné $I$ :
\[
  (X_{i})_{i \in I}, \qquad X_{i} \subset X,
\]
soumise à deux conditions :

\begin{itemize}
\item a) les $X_{i}$ sont fermés et la famille est localement finie ;
\item b) $i \leq j \implies X_{i} \subset X_{j}$.
\end{itemize}

Aucune hypothèse de trivialité locale n'est faite, et les conditions de connexité
et de locale connexité qu'exigera la faisceautisation sur $X$ sont
délibérément différées.

\subsection*{Drapeaux}

On pose
\[
  X_{\Delta_{0}} = \coprod_{i \in I} X_{i}
  = \bigl\{ (x,i) \in X \times I \ \big|\ x \in X_{i} \bigr\},
\]
et la condition a) dit exactement que le morphisme canonique
$X_{\Delta_{0}} \to X$ est \emph{fini} — propre, séparé, à fibres finies ; c'est
aussi une immersion locale. On introduit ensuite la partie fermée
\[
  X_{\Delta_{1}} = \coprod_{i \leq j} \Gamma_{i,j}
  \;\subset\; X_{\Delta_{0}} \times X_{\Delta_{0}},
  \qquad \Gamma_{i,j} \subset X_{i} \times X_{j}
  \text{ le graphe de } X_{i} \hookrightarrow X_{j},
\]
qui n'est autre que le graphe d'une relation d'ordre sur $X_{\Delta_{0}}$ :
\[
  (x,i) \leq (y,j) \iff x = y \text{ et } i \leq j,
\]
soit l'ordre induit par l'ordre produit de $X \times I$, $X$ étant muni de
l'ordre discret. Le couple $(X_{\Delta_{0}}, X_{\Delta_{1}})$ est donc un
\emph{objet ordonné de la catégorie des espaces topologiques} : une catégorie
interne, d'objet des objets $X_{\Delta_{0}}$, d'objet des flèches
$X_{\Delta_{1}}$, de source $\sigma$, de but $b$, et de composition
\[
  X_{\Delta_{2}} = (X_{\Delta_{1}},b) \times_{X_{\Delta_{0}}}
  (X_{\Delta_{1}},\sigma) \longrightarrow X_{\Delta_{1}},
  \qquad (x,y,z) \mapsto (x,z).
\]

Les deux projections ne jouent pas le même rôle, et la dissymétrie est le fait
technique central de la section :

\begin{itemize}
\item $\sigma \colon X_{\Delta_{1}} \to X_{\Delta_{0}}$ est un revêtement étale,
et même constant au-dessus de chaque $X_{i}$, de fibre
$I_{i} = \{ j \in I \mid j \geq i \}$ — donc fini si les $I_{i}$ le sont ;
\item $b \colon X_{\Delta_{1}} \to X_{\Delta_{0}}$ n'est qu'une immersion locale.
\end{itemize}

En degré quelconque, les produits fibrés itérés donnent les drapeaux :
\[
  \begin{aligned}
    X_{\Delta_{r}}
      &= \bigl\{ (x_{0},\ldots,x_{r}) \in X_{\Delta_{0}}^{\,r+1}
        \ \big|\ x_{0} \leq \cdots \leq x_{r} \bigr\} \\
      &\simeq \bigl\{ (x, i_{0},\ldots,i_{r}) \ \big|\
        x \in X_{i_{0}},\ i_{0} \leq \cdots \leq i_{r} \bigr\} \\
      &\simeq \coprod_{i_{*} \in I(\Delta_{r})} X_{i_{0}},
  \end{aligned}
\]
où $I(\Delta_{r}) = \mathrm{Hom}_{\text{ens. ord.}}(\Delta_{r}, I)$ est
l'ensemble des drapeaux de longueur $r$ de $I$. Autrement dit :
$X_{\Delta_{*}}$ est le \emph{nerf} de l'objet ordonné, et il se plonge dans le
produit $X \times I_{\Delta_{*}}$, dont il hérite sa structure simpliciale :
une application croissante $\alpha \colon \Delta_{r'} \to \Delta_{r}$ induit
\[
  \alpha^{*} \colon X_{\Delta_{r}} \longrightarrow X_{\Delta_{r'}},
  \qquad \alpha^{*}(x,i_{*}) = (x,\ i_{*} \circ \alpha).
\]

Deux remarques du manuscrit méritent d'être conservées. D'abord, tout se
généralise en remplaçant l'ensemble ordonné $I$ par une catégorie
$\mathcal{C}$ et la famille par un foncteur de $\mathcal{C}$ vers les parties
fermées de $X$. Ensuite, si l'on part d'une famille sans ordre, on peut toujours
poser $i \leq j \iff X_{i} \subseteq X_{j}$ ; mais la présentation adoptée
n'exige ni que $i \mapsto X_{i}$ soit injective, ni que ce soit un plongement
d'ensembles ordonnés — la réciproque de b) n'est pas requise. C'est une liberté
délibérée, et la section~4 dira pourquoi : elle est ce qui permet de prendre des
images inverses sans avoir à jeter les indices devenus vides.

\subsection*{Ce qui est étale, ce qui ne l'est pas}

La dissymétrie entre $\sigma$ et $b$ se propage à tous les degrés. Si
$\alpha \colon \Delta_{r'} \to \Delta_{r}$ vérifie $\alpha(0) = 0$, alors
$\alpha^{*}$ est un morphisme de revêtements étales au-dessus de $X$ ; en
particulier $X_{\Delta_{r}} \to X_{\Delta_{0}}$, $(x,i_{0},\ldots,i_{r}) \mapsto
(x,i_{0})$, en est un. Plus précisément, au-dessus du composant $X_{i'_{*}}$ de
$X_{\Delta_{r'}}$, le revêtement est constant de fibre l'ensemble des drapeaux
de type $\Delta_{r}$ qui s'envoient sur $i'_{*}$ :
\[
  \alpha^{*-1}\bigl(X_{i'_{*}}\bigr) \;\simeq\;
  X_{i'_{*}} \times \bigl\{ i_{*} \in I(\Delta_{r})
    \ \big|\ i_{*} \circ \alpha = i'_{*} \bigr\}.
\]
Sans la condition $\alpha(0) = 0$, on ne peut affirmer que des immersions
locales — et de même $X_{\Delta_{r}} \to X$ n'est jamais qu'une immersion
locale.\footnote{La raison de la condition $\alpha(0) = 0$ est que le composant
d'un drapeau ne dépend que de son \emph{premier} indice, $X_{i_{*}} = X_{i_{0}}$ ;
une application croissante qui préserve le $0$ ne touche donc pas à l'espace,
seulement à la combinatoire, et c'est ce qui rend $\alpha^{*}$ étale.}

Le manuscrit note enfin qu'on aurait intérêt à éliminer les redondances en se
bornant aux drapeaux \emph{non dégénérés}, c'est-à-dire aux suites strictement
croissantes, ce qui oblige à ne garder que les $\alpha$ strictement croissantes.
C'est la variante semi-simpliciale, et c'est ce que l'auteur appelle
« $\tfrac{1}{2}$ simplicial ».\footnote{Le manuscrit écrit tantôt
« $\tfrac{1}{2}$ simplicial », tantôt « semi-simplicial », pour le même objet.
L'usage d'aujourd'hui distingue : \emph{simplicial} pour le système avec toutes
les applications croissantes, \emph{semi-simplicial} pour la variante à
applications strictement croissantes, qui est celle de son NB de la page 50. Les
deux sont ici en présence, et le texte passe de l'une à l'autre sans le dire.}

\pagerange{50}{53}
\section*{Reconstituer $X$}

La question posée est celle de savoir si le diagramme

\begin{tikzcd}[column sep=large, row sep=small]
& X_{\Delta_{1}} \arrow[dl, "\sigma"'] \arrow[dr, "b"] & \\
X_{\Delta_{0}} \arrow[dr, dashed] & & X_{\Delta_{0}} \arrow[dl, dashed] \\
& X &
\end{tikzcd}

est cocartésien, c'est-à-dire si $X \simeq X_{\Delta_{0}}
\amalg_{X_{\Delta_{1}}} X_{\Delta_{0}}$. Il y faut deux conditions de plus :

\begin{itemize}
\item c) $X = \bigcup_{i} X_{i}$ ;
\item d) pour tous $i,j$, \ $X_{i} \cap X_{j} = \bigcup_{k \leq i,\ k \leq j} X_{k}$.
\end{itemize}

\textbf{Proposition.} Sous a), b), c), d), l'application canonique
\[
  \varinjlim X_{\Delta_{*}} \xrightarrow{\ \sim\ } X
\]
est un homéomorphisme — et l'énoncé vaut pour les deux variantes, drapeaux
quelconques ou drapeaux non dégénérés.

Le corollaire donne la réciproque, et c'est lui qui fait de la construction une
équivalence de données. Soit $(X_{\Delta_{0}}, X_{\Delta_{1}})$ un espace
topologique ordonné tel que $X_{\Delta_{1}}$ soit fermé dans le produit et que,
pour tous $i,j$, la trace $\Gamma_{i,j} = X_{\Delta_{1}} \cap (X_{i} \times
X_{j})$ soit vide ou le graphe d'une immersion fermée $X_{i} \hookrightarrow
X_{j}$. Alors $\sigma$ est un revêtement étale constant, $b$ une immersion
locale, la relation $i \leq j \iff \Gamma_{i,j} \neq \emptyset$ est un ordre sur
$I$, et les $X_{i}$ forment un système inductif d'espaces topologiques à
morphismes de transition des immersions fermées. Si de plus $I$ est fini et
stable par bornes inférieures,
\[
  X = \varinjlim_{i \in I} X_{i} \;\simeq\;
  X_{\Delta_{0}} / \bigl(X_{\Delta_{1}} \rightrightarrows\bigr),
\]
les $X_{i} \hookrightarrow X$ étant des immersions fermées.

\subsection*{Une limite inductive stricte, et pourquoi elle ne suffira pas}

Cette proposition est vraie et n'est pas ce qu'il faudra. La limite inductive
qu'elle exhibe est une limite d'\emph{espaces} : le quotient de
$\coprod_{i} X_{i}$ par la relation d'équivalence engendrée par les inclusions,
c'est-à-dire $\varinjlim_{i \in I} X_{i}$. Elle redonne $X$ comme espace, ce qui
est exact ; mais elle ne dit rien de son type d'homotopie, parce que les
$X_{i}$ sont \emph{fermés} les uns dans les autres. Entre deux morceaux voisins,
la seule chose que ce diagramme enregistre est un bord — et un bord ne dit pas
comment l'un s'approche de l'autre, comme la sphère l'a déjà montré à la section
précédente.

C'est exactement pour cela que la section~3 remplacera les fermés $X_{i}$ par
des \emph{tubes}, qui sont des voisinages. La limite inductive stricte restera
la même ; la limite inductive homotopique, elle, deviendra la bonne.\footnote{La
distinction n'est pas faite sur la page, et le manuscrit n'a pas de mot pour
elle : il écrit $\varinjlim$ dans les deux cas, et parle de « somme amalgamée des
types d'homotopie » quand il veut le second. La théorie qui les sépare — les
colimites homotopiques — lui est contemporaine ; il la nomme lui-même, page 43,
comme la chose qui lui manque et qu'il tient pour standard.}

\subsection*{Faut-il que les intersections soient des morceaux ?}

La condition d) demande que $X_{i} \cap X_{j}$ soit \emph{réunion} des $X_{k}$
en dessous. Faut-il exiger davantage — que la borne inférieure $i \wedge j$
existe et que $X_{i \wedge j} = X_{i} \cap X_{j}$ ?

Le manuscrit répond non, et son contre-exemple est un dessin. Prenez le cercle,
coupé en deux arcs $A$ et $B$ par deux points $a$ et $b$ : quatre strates,
$A \cap B = \{a\} \cup \{b\}$, deux strates connexes dont l'intersection ne l'est
pas. Prenez-en le cône, de sommet $o$ : une stratification locale en $o$ à cinq
strates. « Faut-il l'éliminer pour autant ? Il ne semble pas. La suite le
montrera bien ! » — et il a raison, car ce sont précisément les stratifications
de ce type qui se présentent dans la nature.

\pagerange{53}{55}
\section*{Les strates}

Les $X_{i}$ ne sont pas les morceaux : ils sont les fermés dont les morceaux sont
les différences. Pour tout $x \in X$, l'ensemble $I_{x} = \{ i \in I \mid x \in
X_{i}\}$ est fini par a), et filtrant décroissant par d) ; il a donc un plus
petit élément
\[
  i(x) = \text{le plus petit } i \in I \text{ tel que } x \in X_{i}.
\]
Et pour tout $i$ on pose
\[
  \dot{X}_{i} = \bigcup_{j < i} X_{j} \quad (\text{fermé de } X_{i}),
  \qquad
  X_{i}^{*} = X_{i} \smallsetminus \dot{X}_{i} \quad (\text{ouvert de } X_{i}).
\]

\textbf{Proposition.} Les $X_{i}^{*}$ sont deux à deux disjoints, leur réunion
est $X$, et $i(x)$ est l'unique indice $i$ tel que $x \in X_{i}^{*}$.

\emph{Démonstration.} Soient $i \neq j$ et $x \in X_{i}^{*} \cap X_{j}^{*}$.
Alors $x \in X_{i} \cap X_{j}$, donc par d) il existe $k \leq i,j$ avec
$x \in X_{k}$. Si $k = i$, alors $i \leq j$, donc $i < j$, et $x \in \dot{X}_{j}$,
ce qui contredit $x \in X_{j}^{*}$. Donc $k < i$, d'où $X_{k} \subset \dot{X}_{i}$
et $x \in \dot{X}_{i}$, ce qui contredit $x \in X_{i}^{*}$. La réunion est $X$ et
l'unicité est immédiate par définition de $i(x)$.\footnote{Le signe entre
$\dot{X}_{i}$ et $X_{k}$ est tracé d'un seul trait oblique sur la page et la
transcription hésite ; c'est l'inclusion $X_{k} \subset \dot{X}_{i}$ que
l'argument demande, et c'est elle qui est écrite ici.}

La conséquence à retenir, et que le manuscrit n'énonce pas sous cette forme,
est que $x \mapsto i(x)$ est une application \emph{continue} de $X$ vers $I$
muni de la topologie d'Alexandrov : les parties ouvertes de $I$ y sont les
parties croissantes, et l'image inverse d'une partie décroissante
$\{ j \mid j \leq i \}$ est précisément $X_{i}$, qui est fermé. C'est la
définition moderne d'un espace stratifié par un ensemble ordonné, et elle est ici
entièrement contenue dans les conditions a) à d).\footnote{Cette formulation —
une stratification est un morphisme continu vers un poset alexandrovien — est
celle qui a cours depuis Lurie ; elle est postérieure de trente ans à ces pages.
Ce qui est ici n'est pas la définition mais son contenu : la partition en
strates, l'existence de $i(x)$, et le fait que les $X_{i}$ sont les images
inverses des parties décroissantes. La section~4 ajoutera la moitié qui manque,
en montrant que $I' \mapsto X_{I'}$ commute aux réunions et aux intersections.}

Les mêmes définitions se propagent aux drapeaux :
\[
  \dot{X}_{\Delta_{0}} = \coprod_{i} \dot{X}_{i}, \qquad
  X^{*}_{\Delta_{0}} = X_{\Delta_{0}} \smallsetminus \dot{X}_{\Delta_{0}}
  = \coprod_{i} X_{i}^{*},
\]
et en degré $r$ par image inverse le long de
$\sigma_{r} \colon X_{\Delta_{r}} \to X_{\Delta_{0}}$, ce qui fait de
$\dot{X}_{\Delta_{r}}$ un fermé et de $X^{*}_{\Delta_{r}}$ un ouvert de
$X_{\Delta_{r}}$. Les applications $\alpha^{*}$ avec $\alpha(0) = 0$ respectent
la partition. En revanche — et l'avertissement est de l'auteur — les
$\dot{X}_{\Delta_{r}}$ ne forment pas un espace semi-simplicial : seule la moitié
des applications les préserve.

« Moralement », dit-il, $X^{*}_{\Delta_{r}}$ est la partie non singulière de
$X_{\Delta_{r}}$ et $\dot{X}_{\Delta_{r}}$ son bord. Le propos reste, dans
l'esprit de la section~1, de reconstituer $X$ à partir des strates $X_{i}^{*}$ —
la section~1 étant le cas où $I$ a deux éléments, $X_{0} = Y$ et $X_{1} = X$,
avec $Y = Y^{*}$ et $X^{*} = X \smallsetminus Y$, et où il avait fallu
introduire $\mathcal{V}_{Y,X}$.

\pagerange{56}{60}
\section*{3. Stratifications globales : introduction des tubes}

C'est le mot « il avait fallu » qui commande la section suivante. Pour tout
couple $i \leq j$ on pose
\[
  \overline{\mathcal{V}}_{i,j} = \mathcal{V}_{X_{i},X_{j}}
  \qquad (\text{voisinage tubulaire de } X_{i} \text{ dans } X_{j}),
\]
au sens laissé programmatique par la section~1, et l'on découpe dedans. Le
découpage se fait avec une seule partie fermée auxiliaire :
\[
  R_{i,j} = \bigcup \bigl\{ X_{k} \ \big|\ k < j,\ X_{k} \not\supset X_{i} \bigr\},
\]
la réunion des morceaux de $X_{j}$ qui ne contiennent pas $X_{i}$ — ceux qui
« ne voient pas » la strate depuis laquelle on
regarde.\footnote{Le manuscrit écrit cette partie de deux façons, $R^{*}_{i,j} =
\bigl(\bigcup_{X_{k} \not\supset X_{i}} X_{k}\bigr) \cap X_{j}$ et
$R_{i,j} = \bigcup_{k<j,\ X_{k} \cap X_{i} \neq X_{i}} X_{k}$, et emploie parfois
$R_{i}$ sans second indice pour la réunion non tronquée, qu'il note fermée dans
$X$. Les trois écritures désignent la même chose à la troncation près, et nous
n'en gardons qu'une.}

Le fait qui rend le découpage utilisable est le suivant, et il est démontré sur
la page :
\[
  R_{i,j} \cap X_{i} = \dot{X}_{i}.
\]
En effet, si $X_{k} \not\supset X_{i}$, l'intersection $X_{k} \cap X_{i}$ est,
par la condition d), la réunion des $X_{\ell}$ avec $\ell \leq k,i$ ; chacun de
ces $X_{\ell}$ est strictement contenu dans $X_{i}$, donc $\ell < i$ et
$X_{\ell} \subset \dot{X}_{i}$. L'inclusion inverse est triviale.

\subsection*{Trois espaces, et ce qu'ils valent}

On pose alors
\[
  \mathcal{V}_{i,j} = \overline{\mathcal{V}}_{i,j} \smallsetminus R_{i,j},
  \qquad
  \mathcal{V}^{*}_{i,j} = \overline{\mathcal{V}}_{i,j} \smallsetminus \dot{X}_{j}
  = \mathcal{V}_{i,j} \cap X_{j}^{*},
\]
la seconde égalité valant parce que $R_{i,j} \subset \dot{X}_{j}$. Et
l'identité ci-dessus donne
\[
  \mathcal{V}_{i,j} \cap X_{i} = X_{i}^{*} :
\]
$\mathcal{V}_{i,j}$ est un espace infinitésimal autour de la \emph{strate}
$X_{i}^{*}$, et non autour du fermé $X_{i}$. C'est le gain de tout le découpage.
Pour $i = j$ tout se dégénère proprement :
\[
  \overline{\mathcal{V}}_{i,i} = X_{i}, \qquad
  \mathcal{V}_{i,i} = \mathcal{V}^{*}_{i,i} = X_{i}^{*},
\]
de sorte que les strates elles-mêmes sont des cas particuliers des tubes.

Les pierres de construction sont donc les $\mathcal{V}_{i,j}$ et
$\mathcal{V}^{*}_{i,j}$, avec les inclusions incidentes :

\begin{tikzcd}[column sep=large, row sep=small]
& X_{i}^{*} \arrow[d, "\text{âme}"] \\
\mathcal{V}^{*}_{i,j} \arrow[r] \arrow[d] & \mathcal{V}_{i,j} \\
X_{j}^{*} &
\end{tikzcd}

$X_{i}^{*}$ étant l'\emph{âme} du tube : le mot est de l'auteur, et il est le
bon, puisque c'est la section sur laquelle le tube se rétracte.

\subsection*{La géométrie du tube}

Sous des hypothèses de lissité et de trivialité normale locale — pour une variété
algébrique complexe stratifiée par des sous-variétés fermées, par exemple — les
tubes prennent une forme précise, et c'est cette forme qui justifie qu'on les ait
préférés aux fermés. Le tube épointé $\mathcal{V}^{*}_{i,j}$ est un fibré
localement trivial sur $X_{i}^{*}$, de fibre un cylindre ouvert
$\mathbb{R} \times C_{i,j}$ où $C_{i,j}$ est une variété topologique compacte ;
et $\mathcal{V}_{i,j}$ est le fibré en cônes correspondant, de fibre le cône sur
$C_{i,j}$. Quand $j$ n'est pas un successeur de $i$, $C_{i,j}$ est encore une
variété, mais non compacte.\footnote{$C_{i,j}$ est ce qu'on appelle aujourd'hui
le \emph{lien} de la strate $X_{i}^{*}$ dans $X_{j}$, et la description « voisinage
$=$ fibré en cônes sur le lien » est la définition d'une stratification
\emph{conique}. La chose était disponible : les stratifications de
Thom (1969) et Mather (1970) ont exactement cette propriété, et Siebenmann (1972)
en a fait une définition. Le manuscrit ne cite personne et présente la propriété
comme une conséquence attendue de la trivialité normale locale.}

L'auteur juge alors ses pierres mal taillées : trop générales d'un côté, pas
assez de l'autre pour contenir les morceaux dont il aurait besoin, et il renvoie
à un autre de ses textes.\footnote{Le renvoi est à « Str$_{\infty}$, p.~84 »,
avec une formule $\mathcal{V}^{\,d'_{*}}_{d_{*},d''_{*}} \simeq
\widetilde{N}_{d_{*};(d'_{0},d''_{0})}|\,D^{*}_{d_{*}}$ dont aucun des symboles
n'est défini dans ce dossier. Nous n'identifions pas le texte visé.} La
restriction qu'il retient est de se borner aux couples $i \leq j$ où $j$ est un
\emph{successeur} de $i$ — c'est-à-dire où il n'existe aucun $k$ avec
$i < k < j$. Ce sont ces pierres-là, dites élémentaires, dont le lot suivant
fera le diagramme $\widetilde{I}$.

\subsection*{Les tubes mixtes}

Reste une famille intermédiaire, qui interpole entre $\mathcal{V}^{*}_{i,k}$ et
$\mathcal{V}_{i,k}$. Étant donné $i \leq j \leq k$, on pose
\[
  S_{j,k} = \bigcup \bigl\{ X_{\ell} \ \big|\ \ell < k,\ X_{\ell} \supset X_{j} \bigr\},
  \qquad
  \mathcal{V}^{\,j}_{i,k}
  = \overline{\mathcal{V}}_{i,k} \smallsetminus \bigl( R_{i,k} \cup S_{j,k} \bigr).
\]
Les deux parties $R_{j,k}$ et $S_{j,k}$ répondent à la partition de
$\{\ell \in I \mid \ell < k\}$ selon que $X_{\ell}$ contient ou non $X_{j}$, de
sorte que $R_{j,k} \cup S_{j,k} = \dot{X}_{k}$. Comme $i \leq j$ entraîne
$R_{i,k} \subset R_{j,k}$, on a bien $R_{i,k} \cup S_{j,k} \subset \dot{X}_{k}$,
et la famille est croissante en $j$ :
\[
  \mathcal{V}^{*}_{i,k} \;\subset\; \mathcal{V}^{\,j}_{i,k}
  \;\subset\; \mathcal{V}_{i,k},
  \qquad
  \mathcal{V}^{\,i}_{i,k} = \mathcal{V}^{*}_{i,k}, \qquad
  \mathcal{V}^{\,k}_{i,k} = \mathcal{V}_{i,k},
\]
les deux valeurs extrêmes venant de $R_{i,k} \cup S_{i,k} = \dot{X}_{k}$ d'une
part, de $S_{k,k} = \emptyset$ de l'autre. On retrouve pour $i = j = k$ la strate
$X_{i}^{*}$.\footnote{La page écrit $R_{j,k}$ là où nous écrivons $R_{i,k}$, et
c'est le seul endroit de la section où nous nous écartons d'elle. Avec
$R_{j,k}$, la réunion $R_{j,k} \cup S_{j,k}$ vaut $\dot{X}_{k}$ pour \emph{tout}
$j$, la famille $\mathcal{V}^{\,j}_{i,k}$ est constante et égale à
$\mathcal{V}^{*}_{i,k}$, et les deux identités que la page énonce juste en
dessous sont fausses. Avec $R_{i,k}$, elles sont vraies toutes les deux, la
famille interpole comme annoncé, et la note de marge — $S_{j,k} = \emptyset$ si
$j = k$, $S_{j,k} = X_{j}$ si $k$ est un successeur de $j$ — est exacte. La
correction est donc dictée par les énoncés que la page elle-même en tire. Le
manuscrit ajoute que $R_{j,k} \cap S_{j,k} = \dot{X}_{i}$ ; le membre de gauche
ne dépend pas de $i$, donc l'égalité ne peut se lire telle quelle, et nous ne
l'affirmons pas — ce que l'identité démontrée plus haut donne, c'est
$R_{j,k} \cap X_{j} = \dot{X}_{j}$.}

Le tube mixte $\mathcal{V}^{\,j}_{i,k}$ tend lui aussi à être un fibré localement
trivial de base $X_{i}^{*}$, mais sa fibre peut être singulière : dans le cas
$\mathcal{V}^{\,k}_{i,k} = \mathcal{V}_{i,k}$ avec $k$ successeur de $i$, la
singularité est isolée ; en général elle peut être de type arbitraire — il suffit
de prendre $X_{i} = X_{i}^{*}$ réduit à un point pour la mettre tout entière dans
la fibre. C'est ce qui diminue l'intérêt des tubes mixtes comme pierres, et
c'est pourquoi le programme annoncé à la fin de la section est celui-ci :
exprimer tout avec les $X_{i}^{*}$ et les tubes élémentaires
$\mathcal{V}_{i,j}$, $\mathcal{V}^{*}_{i,j}$ ($j$ successeur de $i$) ; voir
ensuite si certaines expressions se simplifient en admettant tous les
$i \leq j$ ; et seulement en dernier lieu, en vue du cas des diviseurs à
croisements normaux, voir comment les $\mathcal{V}^{\,j}_{i,k}$ s'insèrent dans
ce formalisme.

C'est ce programme que le dernier lot exécute, dans l'ordre annoncé, et qu'il
laisse inachevé à sa troisième étape.

\end{document}
